\section{Introduction}\label{sec:intro}

Organizations allocate their users across channels of different quality and cost: a specialized service center versus a general one, a hospital versus primary care, a personal advisor versus an app. Before reallocating, they want to know how much would be lost by moving a group to the basic channel. A randomized experiment would answer this, but it is often expensive, slow or infeasible; what is available are observational data from each channel. The standard approach reweights the units that the basic channel serves today so that they resemble the target group, with weights estimated by a classifier. Formally, this is a change of measure.

This paper is about a feature of channel data that the standard approach ignores: \emph{each channel records covariates in its own way}. A personal advisor writes down a commercial history that the app never stores; a hospital runs tests that primary care does not; one channel records what customers declare, while another measures it. We show that this breaks reweighting in two ways, one visible and one invisible.
\begin{itemize}
  \item \textbf{A variable recorded by only one channel} breaks support by construction, however comparable the units are: the unsupported mass is at least the fraction of the target group with data, and the ideal classifier reveals it by being certain ($p=1$) on those profiles (Proposition~\ref{prop:breaks}); in an experimental benchmark (NSW versus CPS), a top-coded earnings variable breaks support at its edge and the estimated classifiers miss it (Section~\ref{sec:lalonde}). Dropping the variable restores identification when it does not predict the basic-channel outcome given the other covariates (Proposition~\ref{prop:drop}); if it does, neither keeping nor dropping it identifies the estimand.
  \item \textbf{A variable measured with a different instrument in each channel} can preserve support. The estimator then converges to a shifted value, which we derive exactly, and which no analysis of the observed data can distinguish from the truth (Theorem~\ref{thm:measurement}). With public data in which both instruments exist for the same people (self-reported and measured body mass index in NHANES), the derived shift approximately matches the observed bias in 14 of 15 registered checks (the exception is explained by the outcome model), and the usual overlap diagnostics improve, relative to correct measurement, rather than flag it (Section~\ref{sec:nhanes}).
\end{itemize}

To state these results we need the reweighting machinery in measure-theoretic form (Section~\ref{sec:identification}). Most of it is known, and we do not claim it: the estimand is the counterfactual mean that enters an average treatment effect on the treated, the odds weights $p/(1-p)$ are those of \citet{hirano2003}, estimating density ratios with classifiers is standard \citep{qin1998,bickel2009,sugiyama2012}, and the doubly robust theory with cross-fitting is that of \citet{chernozhukov2018}. Two readings are new, and Sections~\ref{sec:channel} and~\ref{sec:instruments} rely on them. The unsupported part of the target law is exactly its restriction to $\{p=1\}$, which is von Neumann's construction of the Lebesgue decomposition \citep[Theorem 6.10]{rudin1987} read through the classifier; and without support the doubly robust estimator silently fills that part by extrapolating the outcome model (Theorem~\ref{thm:lebesgue} and Corollary~\ref{cor:drnosupport}).

We also give tools for practice. When the target group is a cutoff on a known score, the weights are the ordinary channel propensity with the cutoff imposed (Proposition~\ref{prop:cutoff}), as in subgroup analysis \citep{rassen2012}; estimating that propensity with all units of the premium channel reduced the dispersion by up to 59\,\% in our simulation when overlap was good, but not when it was poor. We give the exact bias of capping the true weights with normalization (Proposition~\ref{prop:trimming}), which complements \citet{ma2020}, and adapt the bounds of the marginal sensitivity model \citep{tan2006,dorn2023} to this estimand, with a calibration through the channel model (Section~\ref{sec:sensitivity}). Besides the two public-data analyses, a simulation tests the other results: for each scenario we computed the limit of each estimator, and registered further predictions, before running it, and we report the predictions that failed (Sections~\ref{sec:design} and~\ref{sec:results}). Trimming with estimated weights and the behavior of the bootstrap are the subject of a companion note in preparation; here we only summarize the results.

Section~\ref{sec:setup} sets up the problem and Section~\ref{sec:identification} the reweighting machinery. Sections~\ref{sec:channel} and~\ref{sec:instruments} contain the two failures, each followed by public data: the visible one with the NSW/CPS benchmark (Section~\ref{sec:lalonde}), the invisible one with NHANES (Section~\ref{sec:nhanes}). Sections~\ref{sec:cutoff}--\ref{sec:sensitivity} collect tools for practice that do not depend on the measurement problem: known-cutoff weights, trimming and inference, and sensitivity bounds. Sections~\ref{sec:design} and~\ref{sec:results} report the simulation, and Section~\ref{sec:discussion} concludes.

We use the language of measure theory (Radon--Nikodym derivatives) rather than density ratios. This is not a contribution but the standard formulation; it handles mixed covariates (continuous, categorical and with missing values) without assuming that densities exist.

\paragraph{Related work.} The same estimand, odds weights and doubly robust influence function appear in the literature on generalizing and transporting trial results to a target population \citep{westreich2017,dahabreh2019}. \citet{heckman1997} found empirically that matching removes much of the bias when the comparison group comes from the same market and answers the same questionnaire as participants; Theorem~\ref{thm:measurement} formalizes that lesson as an indistinguishability result. \citet{cook2008} and \citet{steiner2011} also document the importance of measuring covariates in the same way and reliably, and that differential measurement error cannot be corrected without external information is known in the measurement-error literature; \citet{imai2010} study nonparametric identification and sensitivity analysis under differential measurement error of the treatment. On partially observed covariates, the missingness-pattern approach \citep{dagostino2000,blake2020} requires the variable not to confound where it is missing; our channel variables are the case in which the treatment fixes the pattern, and the condition we use is an outcome-based covariate selection criterion \citep{deluna2011,vanderweele2011}. On diagnostics, \citet{westreich2011} show that a high AUC of the propensity model is neither necessary nor sufficient to control confounding; our simulation adds that it does not measure overlap either. For diagnosing positivity violations, see \citet{petersen2012} and \citet{damour2021}. For limited overlap and trimming, see \citet{crump2009}, \citet{ma2020,ma2023}, \citet{liu2023} and \citet{manamkoong2025}, whose approach of extrapolating the part without overlap under a smoothness assumption complements Corollary~\ref{cor:drnosupport}; we discuss them where they arise (Section~\ref{sec:inference}).
